% Source: https://artofproblemsolving.com/community/c3h1517876_geometry
% Stable identity: aops-topic-1517876-q01
For a celebration, an invitation was designed by overlapping three squares, as shown. Each square has a positive integer side length (in centimetres). Side $AB$ of the smallest square lies along side $AC$ of the middle square, which lies along side $AD$ of the largest square; the three squares share the corner $A$ and lie on the same side of line $AD$.

\begin{asy}
usepackage("amsmath");
size(4.5cm);
defaultpen(fontsize(10pt));
real a=3,bb=5.5,c=8;
draw((0,0)--(c,0)--(c,c)--(0,c)--cycle,linewidth(1));
draw((bb,0)--(bb,bb)--(0,bb),linewidth(1));
draw((a,0)--(a,a)--(0,a),linewidth(1));
label("$A$",(0,0),SW);label("$B$",(a,0),S);label("$C$",(bb,0),S);label("$D$",(c,0),SE);
label("\mbox{}",(min(currentpicture,user=true)+max(currentpicture,user=true))/2,black); // keeps the text after the figure black when LaTeX typesets the labels inline
\end{asy}

The area of the middle square not covered by the smallest square is $33\text{ cm}^2$. If $BC=CD$, determine all possible side lengths of the three squares. (The figure is not to scale.)
